\documentclass[10pt,a4paper]{amsart}
\usepackage[margin=19mm]{geometry}
\usepackage{amsmath,amssymb,mathtools}
\usepackage[scr=rsfs,cal=euler]{mathalfa}
\usepackage{xfrac}
\usepackage[unicode,hidelinks,bookmarks=false]{hyperref}
\newcommand{\ie}{i.e.\ }
\newcommand{\cf}{cf.\ }
\newcommand{\resp}{respectively\ }
\newcommand{\Subsection}{\subsection}
\providecommand{\nobreakdash}{\nobreak\hbox{-}}
\setlength{\emergencystretch}{2em}
\multlinegap0pt
\usepackage{bm}
\usepackage{bbm}
\usepackage{dsfont}
\usepackage{xspace}
\usepackage{cancel}
\usepackage{mathtools}
\usepackage{amsthm}
\makeatletter
\@ifundefined{Lemma}{\newtheorem{Lemma}{Lemma}}{}
\@ifundefined{Theorem}{\newtheorem{Theorem}{Theorem}}{}
\makeatother
\title[Multicolor vector space Ramsey numbers over the binary field]{Multicolor vector space Ramsey numbers over the binary field}
\author{Anurag Bishnoi, Gaurav Kucheriya}
\date{Source: arXiv:2607.17263v1 (CC BY 4.0)}
\begin{document}
\maketitle

\noindent\emph{Source and licence.} Multicolor vector space Ramsey numbers over the binary field. Version arXiv:2607.17263v1 (CC BY 4.0), distributed under the Creative Commons Attribution 4.0 International licence, as recorded in the arXiv licence metadata for this exact version and shown on the abstract page https://arxiv.org/abs/2607.17263v1. Full rights notice: licenses/arxiv-2607.17263-CC-BY-4.0.txt.\par\smallskip

\noindent\textit{Editorial scope.} Counted unit: the Theorem whose source label is \texttt{None} in the article (source0.tex lines 198--200, source label \texttt{None}), delivered together with the article's own complete original proof of it (source0.tex lines 201--230, one \texttt{proof} environment of 30 lines). No mathematics is written, repaired or re-derived: statement and proof are the article's own text line for line. Required context retained uncounted: lem:simplex (lines 138--143); lem:isomorphism (lines 149--163). Every definition, display and label of those lines is retained unchanged. Omitted from this package and not redistributed: 1--137, 144--148, 164--197, 231--360. Nothing inside a retained excerpt is omitted or altered: every retained body is delivered exactly as the article's lines are written, so no omission receipt of any kind accompanies this package. Citations: the retained excerpts carry no citation command at all, so no bibliography entry of the article is transcribed and no reference is redistributed. Reference adaptation: none was needed and none was made; every reference inside the delivered unit resolves inside it, so no unresolved reference or citation is produced. Printed slips kept literal: no printed word, symbol, hypothesis or number of the counted statement or of its proof was altered. Layout and class: the article's own class is \texttt{amsproc} with packages including amssymb, amsmath, amsthm, lmodern, tikzsymbols, amsaddr, framed, booktabs, float, xcolor, enumerate, mathtools, url, fontenc; the wrapper is the lane's pinned \texttt{amsart} editorial wrapper at 10pt on A4, which restates the theorem-like environments the retained text needs and adds the article's own macro definitions that the retained text needs (none), with extra wrapper packages (bm, bbm, dsfont, xspace, cancel, mathtools). The delivered numbering is the unit's own. Assets: no retained span contains a figure environment, a graphics-inclusion command, a PDF-page inclusion, a table or a TikZ picture; no image file is copied into this package, the package declares no asset dependency and no image file is redistributed with it. Licence: version arXiv:2607.17263v1 of this article is distributed under the Creative Commons Attribution 4.0 International licence, as recorded in the arXiv licence metadata for this exact version and shown on the abstract page https://arxiv.org/abs/2607.17263v1. A full rights notice is delivered at licenses/arxiv-2607.17263-CC-BY-4.0.txt. Characters: every retained excerpt is pure ASCII, so the delivered engine, XeLaTeX with its default Latin Modern text font, reports no missing character.
\begin{Lemma} 
\label{lem:simplex}
    The $q$-ary simplex code $\mathcal{S}_{q,t}$ has parameters 
    \begin{equation*} \left[\frac{q^t-1}{q-1},\,t,\,q^{t-1}\right]_q. \end{equation*} 
    Moreover, every nonzero codeword of $\mathcal{S}_{q,t}$ has Hamming weight $q^{t-1}$. 
\end{Lemma}

\begin{Lemma}\label{lem:isomorphism}
    Let $C \leq \mathbb{F}_q^m$ be a $t$-dimensional vector subspace. 
    Let $x_1, \dots, x_m$ be vectors in $\mathbb{F}_q^n$. 
    Define
    \[
        S \coloneqq \left\{\sum_{i \in [m]} c_i x_i :
        (c_1, \dots, c_m) \in C\right\}.
    \]
    If, for every $(c_1,\dots,c_m)\in C$,
    \[
        \sum_{i \in [m]} c_i x_i = \vec{0}
        \implies c_i = 0 \text{ for all } i,
    \]
    then $S \cong C$.
\end{Lemma}

\begin{Theorem}
    If $\chi_2(t; n) \leq k$, then $R(K_s^{(r)}; k + 1) > 2^n$, where $r = 2^{t - 1}$ and $s = 2^t - 1$.
\end{Theorem}

\begin{proof}
    Let $\chi$ be a proper $k$-coloring of $\mathrm{PG}(n - 1, 2)$, that is, there are no monochromatic $(t - 1)$-flats under $\chi$. 
    Let $r = 2^{t - 1}$ and $s = 2^t - 1$.
    Define a coloring $c$ on the $r$-element subsets of $V = \mathbb{F}_2^n$ by 
    \begin{equation*}
        c(\{x_1, \dots, x_r\}) = 
        \begin{cases}
            \star, & \text{if } x_1 + \cdots + x_r = 0, \\
            \chi(x_1+ \dots+ x_r), & \text{otherwise.}
        \end{cases}
    \end{equation*}
    % $c(\{x_1, \dots, x_r\}) = 0$ if $x_1 + \cdots + x_r = 0$ and $c(\{x_1, \dots, x_r\}) = \chi(x_1 + \cdots + x_r)$ otherwise. 
    We will show that there is no monochromatic $K_s^{(r)}$ in this coloring of the complete $r$-uniform hypergraph on $V$, which proves that $R(K_s^{(r)}; k+1) > 2^n.$
   
    Let $x_1, \dots, x_s$ be a set of distinct vectors in $\mathbb{F}_2^n$.
    First, assume that they induce a monochromatic $K_s^{(r)}$ in color $\star$, that is, 
    $\sum_{i \in T} x_i = 0$ for all subsets $T$ of $[s]$ with $|T| = r$. 
    Since $r+1 \leq s$, we can take $T_1 = \{1, \dots, r\}$ and $T_2 = \{1, \dots, r - 1, r+1\}$ as two such subsets. 
    Then we get $\sum_{i \in T_1} x_i = \sum_{i \in T_2} x_i = 0$, which implies $x_r = x_{r + 1}$, a contradiction to the fact that we have distinct vectors. 

    Now assume that these vectors give a monochromatic $K_s^{(r)}$ in one of the $k$ colors of $\chi$. 
    In particular, $\sum_{i \in T} x_i \neq 0$ for all $T \subseteq [s]$ of size $r$. 
     
    Let $C$ be the binary simplex code $\mathcal{S}_{2,t}$.
    By Lemma~\ref{lem:simplex}, it has parameters $[2^{t} - 1, t, 2^{t-1}]_2$ and every nonzero codeword in $C$ has Hamming weight $2^{t - 1}$.
    Define $S = \left\{\sum_{i \in [s]}c_i x_i : (c_1, \dots, c_s) \in C \setminus \{\vec{0}\}\right\}$ and $\overline{S} = S \cup \{\vec{0}\}$. 
    The Hamming weight of each nonzero codeword in $C$ is $r = 2^{t - 1}$, and the sum of any $r$ vectors among $x_1, \dots, x_s$ is nonzero.
    By Lemma~\ref{lem:isomorphism}, for $q = 2$, we get that $\overline{S}$ is a $t$-dimensional vector subspace of $\mathbb{F}_2^n$, that is $S$ corresponds to a copy of $\mathrm{PG}(t - 1, 2)$ in $\mathrm{PG}(n - 1, 2)$.
    All elements of $S$ receive the same color under $\chi$, which is a contradiction to the fact that $\chi$ is a proper coloring. 
\end{proof}

\end{document}
